\paragraph{Carry the premium-allocation groups into a dated interim book.}
\label{sec:insurance-paa-dated-book}
The short-duration case in Chapter~\ref{ch:insurance} provides another route
from sector measurements to the shared book. It starts with cash and equity
20, receives premium 120, pays reinsurance premium 24, and observes half the
supplied uniform service. Its purpose differs from the loss-component schedule
that follows: both issued and held groups must remain identifiable through
their coverage, claims, service and cash movements.

\code{InsurancePAADatedBook/1} proposes the eight-event sequence. It uses
the unchanged interfaces \code{DeclaredCompanyProgram/2},
\code{CompanyBook/1} and \code{CompanyBookReport/1}.
Its request contains a profile identifier, payload
and payload digest. The payload supplies entity and currency, the explicit
2024 coverage calendar, distinct issued and held groups, their ownership and
link, the selected measurements, eligibility findings, cash classification
and exact source edition. Money and fractions use canonical rational strings.
The full closed fields and serialized examples are included with the companion;
unknown extra fields do not acquire meaning through an informal account label.

The bridge supports three retained measurement cases, a named or renamed
account chart, and an optional unrelated equipment balance of 17 with matching
equity. Group identifiers remain in typed memory and event inputs. Each method
requires the correct date, preceding step and opening balances; the company
engine checks the resulting journal, memory and state before the next event.
The unrelated balances survive. The returned \code{proposed\_conditional}
result contains the complete compiled program, native execution, group records
and unresolved conditions, including every intermediate journal and state.

The final report separately presents all mapped opening and closing accounts,
assets, liabilities, income, expenses, equity movement and gross cash origins.
The base case has cash 101, reinsurance coverage asset 12 and recovery asset 6
against issued coverage 60 and incurred claims 12. Profit is 27 and equity 47.
The transferred-risk variant changes the recovery asset to 6.5; the separate
non-performance variant changes it to 4. Cash remains 101. Both groups remain
gross. The report uses the full coverage-year interval but observes only the
2 July state; it neither supplies the second half-year nor closes the year.

Missing eligibility or measurement findings return \code{unresolved} without
a compiled execution. Contradicted findings or measurements outside the selected
finite domain return \code{unsupported}. Malformed identities, links or digests
return \code{invalid}. A refusal by the compiled native program preserves its
stopped execution and preceding state; the bridge returns \code{native\_refusal}
without inventing a final report. The complete expected states and financial
reports were authored separately before execution. They cover both charts,
the unrelated balance and sequence/date/group failures.

These are proposed conditional entries. Source, group, chronology, accounting
and posting admission remain false, as do durable commitment and fiscal or
note completeness. The supplied uniform risk release, short claims timing,
PAA eligibility and operating cash classification remain assumptions. The
retained AASB~17 edition qualifies the same narrow example as the chapter;
replay does not establish actual applicability. The exact contract, source
qualification, expectations and internal reproduction archive are under
\path{companions/insurance-paa-dated-book/} in this package. The original record
remains under \path{data/accounting-master/insurance-paa-dated-book-20261027/}.

